\chapter{基于双通道注意力的双曲图神经网络}

本章采用第三章提出的隐式层次建模范式，借助双曲空间的几何特性捕获路网层次信息。

\section{问题定义与研究动机}

设路网为有向图 $G = (V, E, \mathbf{X})$，其中 $V = \{v_1, v_2, \ldots, v_N\}$ 为节点集合（每个节点对应一条路段），$E \subseteq V \times V$ 为边集合（表示路段间的拓扑连接），$\mathbf{X} \in \mathbb{R}^{N \times F}$ 为节点属性矩阵（每行对应路段的多维特征向量，包括道路类型、限速、车道数、几何形状等）。

路网表示学习的目标是学习映射 $f: V \to \mathbb{R}^d$，将每个路段节点映射为低维连续向量 $\mathbf{h}_{v_i} \in \mathbb{R}^d$，使得该向量同时编码路段的属性特征和拓扑结构信息，可直接用于下游任务。

本章聚焦于路段类型分类（Road Type Classification）任务：给定路网图 $G$ 和部分路段的类别标签 $\{y_{v_i} \in \mathcal{C} : v_i \in \mathcal{V}_{\text{train}}\}$（$\mathcal{C}$ 为道路类型集合），学习分类器 $g: \mathbb{R}^d \to \mathcal{C}$，以预测测试集路段的类别。道路类型通常包括高速公路（Motorway）、干道（Primary Road）、次要道路（Secondary Road）、居民区道路（Residential）等，反映路网的层次功能组织。

本章考虑两种评估场景。城市内分类（Intra-city Classification）中，训练集和测试集均来自同一城市路网，通过节点掩码划分。跨城市分类（Cross-city Classification）中，训练集来自若干源城市，测试集来自未见过的目标城市，评估模型的跨域泛化能力。跨城市场景对层次结构编码的泛化性提出了更高要求——如第三章分析所述，如果层次关系的编码依赖于几何先验而非特定拓扑细节，则更有望实现跨城市迁移。

如第三章所论述，隐式范式的核心思路是利用双曲空间的度量性质隐式捕获层次信息。本章选择庞加莱球模型，原因在于分类任务更关注路段的判别性特征，庞加莱球的共形性有助于注意力权重的计算，且在该模型中双曲距离可以直接作为注意力系数的相似性度量。

\section{模型总体框架}

本节介绍用于学习路网表示的 HRGNN。首先对模型架构进行总体概述，然后详细描述各关键组件。

\subsection{模型总览}

HRGNN 旨在利用双曲几何 \cite{Krioukov2010} 建模层级结构的优势，学习路网的有效表示。如图~\ref{fig:architecture} 所示，模型由以下主要组件构成：

\begin{enumerate}
\item \textbf{路网图构建（Road Graph Construction）：} 通过将路段转化为节点、将其连接转化为边，对原始路网图进行增强特征表示的变换。

\item \textbf{双曲邻居聚合（Hyperbolic Neighbor Aggregation）：} 这一核心组件通过四个相互关联的子模块，以流水线方式处理路网特征。首先，特征投影与变换（Feature Projection and Transformation）模块利用指数映射将输入路段特征映射至双曲空间，然后施加双曲线性变换以生成潜在表示。其次，这些双曲嵌入被送入双路径注意力（Dual-Pathway Attention）模块，该模块在双曲空间中计算注意力机制，同时融合局部邻居采样和基于 DFS 的拓扑邻居采样，以及基于双曲距离的注意力系数。第三，被加权的表示通过多层聚合（Multi-layer Aggregation）模块，利用自环机制跨多层捕获多尺度结构信息以增强节点表示。最后，自适应融合策略（Adaptive Fusion Strategy）模块利用可学习的融合权重动态组合局部和 DFS 路径的表示，生成最终有效建模路网层级特性的双曲嵌入。

\item \textbf{任务模块与模型训练（Task Module and Model Training）：} 最终双曲嵌入被投影至欧几里得空间，并通过分类头处理以完成下游任务。模型采用结合交叉熵损失与正则化项的复合损失函数，并通过对神经网络参数和双曲曲率的交替优化进行训练。
\end{enumerate}

\begin{figure}[t]
\centering
\includegraphics[width=\textwidth]{figures/figure1.png}
\caption{HRGNN 模型架构总览。模型通过图构建、特征提取和邻居聚合，将路网变换为双曲表示。}
\label{fig:architecture}
\end{figure}

\subsection{路网图构建}

在标准路网图中，交叉口被表示为节点，路段被表示为边。为便于 GNN 对节点特征进行操作，我们采用线图变换（Line Graph Transformation）来在节点上处理路段特征（长度、宽度、限速）\cite{Chai2020}。

变换后的线图将完整的路段属性作为节点特征。每个节点通过标准化几何向量编码几何属性，该向量通过等距采样点捕获路段的空间形态，确保所有路段具有统一的特征维度。功能属性包括表示不同速度类别的独热编码（One-hot Encoding）限速，以及反映从高速公路到地方街道层级组织的道路类型分类。空间上下文通过中点坐标提供位置参考，所有几何特征相对于路段中点进行归一化以确保平移不变性。此外，基于度中心性（Degree Centrality）的连通性度量捕获每个路段在整体网络层级中的拓扑重要性；长度属性则提供补充几何表示的尺度信息。

给定路网图 $G = (V, E)$，其中 $V$ 表示交叉口，$E$ 表示路段，我们构建线图 $L(G) = (V_L, E_L)$，其中每个节点 $v_L \in V_L$ 对应一条路段 $e \in E$。若路段 $e_i, e_j \in E$ 共享同一交叉口，则存在边 $(v_i, v_j) \in E_L$。该变换 \cite{Gharaee2021} 支持直接将路段属性用作节点特征，使其可供 GNN 操作使用，同时保留网络的连通结构。

\subsection{双曲几何}

双曲几何为涉及层级数据结构的机器学习应用提供了强大的数学框架。与欧几里得几何不同，双曲空间具有负曲率和指数体积增长特性，天然适合嵌入社交网络、知识图谱和自然语言处理任务中常见的树状结构。

本节建立数学基础：首先探讨双曲几何的基本属性，重点介绍庞加莱球模型（Poincar\'e Ball Model）；其次介绍旋转向量空间（Gyrovector Spaces）作为在双曲几何中实现向量化操作的代数框架。

\subsubsection{双曲几何基础}

双曲几何是一种非欧几里得几何，它否定欧几里得第五公设（即平行公设）的有效性。在欧几里得几何中，过直线外一点有且仅有一条直线与该直线平行；而在双曲几何中，存在无穷多条这样的平行线 \cite{Sarkar2011,Nickel2017}。

双曲空间具有多个使其特别适合表示层级结构的关键属性 \cite{Sala2018}：双曲空间具有恒定负曲率，导致其体积随半径呈指数级而非多项式增长，这使得层级结构的嵌入更为高效，因为所需空间随与原点距离的增加呈指数增长。双曲空间可以被视为树的连续类比——双曲空间中体积的指数增长对应于树中每层节点数量的指数增长，使双曲几何天然适合嵌入树状结构。双曲空间可以以任意低的失真嵌入树状结构，而欧几里得空间通常需要高维度才能达到可比效果。

双曲几何有多种数学模型，包括庞加莱球模型、洛伦兹（Lorentz）模型（或超球面模型）、克莱因（Klein）模型、半球模型和庞加莱半平面模型。本文主要采用庞加莱球模型，因为它具有共形性和直观的可视化属性。

\subsubsection{旋转向量空间操作}

旋转向量空间为双曲几何提供了代数框架，类似于向量空间为欧几里得几何提供支撑。该框架支持在双曲空间中定义向量化操作 \cite{Ungar2008}。

在旋转向量空间中，向量加法和标量乘法等操作被推广为其双曲对应形式——莫比乌斯加法（M\"obius Addition）和莫比乌斯标量乘法（M\"obius Scalar Multiplication）。这些操作在保持双曲结构的同时，允许执行类似欧几里得空间的计算。

对于庞加莱球模型，定义 $d$ 维球为
\begin{equation}
\mathbb{B}_c^d := \{x \in \mathbb{R}^d : c|x|^2 < 1\}
\label{eq:poincare_ball}
\end{equation}
其中 $c \geq 0$ 为曲率参数。当 $c = 0$ 时，$\mathbb{B}_c^d = \mathbb{R}^d$ 即为欧几里得空间；当 $c > 0$ 时，$\mathbb{B}_c^d$ 为半径为 $1/\sqrt{c}$ 的开球。

两点 $x, y \in \mathbb{B}_c^d$ 的莫比乌斯加法定义为
\begin{equation}
x \oplus_c y := \frac{(1 + 2c\langle x,y\rangle + c|y|^2)x + (1 - c|x|^2)y}{1 + 2c\langle x,y\rangle + c^2|x|^2|y|^2}
\label{eq:mobius_addition}
\end{equation}

点 $x \in \mathbb{B}_c^d \setminus \{0\}$ 与标量 $r \in \mathbb{R}$ 的莫比乌斯标量乘法定义为
\begin{equation}
r \otimes_c x := \frac{1}{\sqrt{c}}\tanh\left(r\tanh^{-1}(\sqrt{c}|x|)\right)\frac{x}{|x|}
\label{eq:mobius_scalar}
\end{equation}

这些操作使我们能够在双曲空间中执行计算，同时保留其几何属性，是开发在该空间中运行的神经网络的关键基础。

\subsection{双曲邻居聚合}

本节介绍双曲邻居聚合模块的设计，该模块利用双曲几何的独特优势，有效捕获局部连通模式和长程依赖关系。

该模块首先将路段特征投影至庞加莱球模型，然后构建能够通过重要性引导的深度优先搜索策略和自适应融合机制同时处理邻居信息和远距离拓扑关系的注意力网络。这一设计不仅保留了双曲空间的几何属性，还通过多层聚合和动态权重调整增强了模型对复杂路网结构的表达能力。完整的流程详见算法1。

\begin{algorithm}[H]
\caption{自适应融合双曲注意力}
\label{alg:hyperbolic_attention_formalized}
\begin{algorithmic}[1]
\REQUIRE 图 $G(V,E)$，节点特征 $h_v$，局部邻域 $N_l$，全局邻域 $N_g$
\ENSURE  更新后的节点表示 $h_{v\_new}$

\vspace{0.5em}
\STATE // 处理局部邻域
\FOR{每个节点 $v \in V$}
    \STATE $d_{vu}^{l} \leftarrow d_{\mathcal{H}}(v, u),\ \forall\, u \in N_l(v)$
        \COMMENT{双曲距离}
    \STATE $\alpha_{vu}^{l} \leftarrow \operatorname{softmax}\!\left(-\sqrt{(d_{vu}^{l})^2}\right),\ \forall\, u \in N_l(v)$
        \COMMENT{局部注意力权重}
    \STATE $h_{v\_local} \leftarrow \displaystyle\sum_{u \in N_l(v)}\alpha_{vu}^{l} \cdot h_u + (1+\varepsilon) \cdot h_v$
        \COMMENT{局部邻域聚合}
\ENDFOR

\vspace{0.5em}
\STATE // 基于动态权重处理全局邻域
\FOR{每个节点 $v \in V$}
    \STATE $z_v \leftarrow f_{\mathrm{extract}}(h_v)$
        \COMMENT{特征提取}
    \STATE $\beta_v \leftarrow \sigma(W_{\beta} \cdot z_v + b_{\beta})$
        \COMMENT{节点自适应权重}
    \STATE $d_{vu}^{g} \leftarrow d_{vu} \cdot (1 - \lambda \cdot \beta_v),\ \forall\, u \in N_g(v)$
        \COMMENT{修正后的全局距离}
    \STATE $\alpha_{vu}^{g} \leftarrow \operatorname{softmax}\!\left(-\sqrt{(d_{vu}^{g})^2}\right),\ \forall\, u \in N_g(v)$
        \COMMENT{全局注意力权重}
    \STATE $h_{v\_global} \leftarrow \displaystyle\sum_{u \in N_g(v)}\alpha_{vu}^{g} \cdot h_u + (1+\delta\varepsilon) \cdot h_v$
        \COMMENT{全局邻域聚合}
\ENDFOR

\vspace{0.5em}
\STATE // 自适应融合局部信息与全局信息
\FOR{每个节点 $v \in V$}
    \IF{$\mathrm{adaptive} = \mathrm{true}$}
        \STATE $\gamma_v \leftarrow \sigma(W_{\gamma} \cdot h_{v\_local} + b_{\gamma})$
            \COMMENT{自适应融合系数}
        \STATE $h_{v\_fused} \leftarrow \gamma_v \cdot h_{v\_local} + (1 - \gamma_v) \cdot h_{v\_global}$
    \ELSE
        \STATE $h_{v\_fused} \leftarrow \dfrac{1}{2}\!\left(h_{v\_local} + h_{v\_global}\right)$
    \ENDIF
    \STATE $h_{v\_new} \leftarrow \phi(W \cdot h_{v\_fused} + b)$
        \COMMENT{MLP 非线性变换}
\ENDFOR

\RETURN $h_{v\_new}$
\end{algorithmic}
\end{algorithm}

\subsubsection{特征投影与变换}

构建线图后，我们将节点特征（即路段属性）投影至双曲空间。这一投影步骤是必要的，因为输入特征可能不自然地存在于双曲空间的约束范围内，而双曲几何为建模层级路网结构提供了独特优势。

我们使用指数映射（Exponential Map）将特征从原点处的切空间投影到庞加莱球 \cite{Mathieu2019}。对于节点 $i$ 的特征向量 $f_i$，投影执行如下：
\begin{equation}
p_i = \exp_0^c(f_i) = \tanh(\sqrt{c}|f_i|)\frac{f_i}{\sqrt{c}|f_i|}
\label{eq:exp_map}
\end{equation}
其中 $c > 0$ 为双曲空间的曲率参数。

投影后，我们施加双曲线性变换以获得更高层的潜在表示。该变换通过莫比乌斯矩阵向量乘法 \cite{Ganea2018} 实现：
\begin{equation}
z_i^H = W \otimes_c x_i^H = \frac{1}{\sqrt{c}}\tanh\left(\frac{|Wx_i^H|}{|x_i^H|}\tanh^{-1}(\sqrt{c}|x_i^H|)\right)\frac{Wx_i^H}{|Wx_i^H|}
\label{eq:mobius_matvec}
\end{equation}
其中 $W \in \mathbb{R}^{D \times F}$ 为可学习权重矩阵，$\otimes_c$ 表示曲率为 $c$ 的双曲空间中的莫比乌斯矩阵向量乘法。该操作在保持双曲结构的同时对输入特征进行变换 \cite{Chami2019}，使模型能够学习保留路网层级性质的表示。

\subsubsection{双曲距离计算}

对于线图中的任意边 $(i,j) \in E$，我们使用庞加莱球模型计算节点间的双曲距离：
\begin{equation}
d_{\mathbb{H}}(i,j) = \frac{2}{\sqrt{c}}\tanh^{-1}\left(\sqrt{c}|z_i \ominus_c z_j|\right)
\label{eq:hyperbolic_distance}
\end{equation}
其中莫比乌斯减法（M\"obius Subtraction）$\ominus_c$ 定义为
\begin{equation}
u \ominus_c v = u \oplus_c (-v) = \frac{(1 + 2c\langle u,-v\rangle + c|v|^2)u + (1 - c|u|^2)(-v)}{1 + 2c\langle u,-v\rangle + c^2|u|^2|v|^2}
\label{eq:mobius_subtraction}
\end{equation}

\subsubsection{双路径双曲注意力}

本文定义了两条用于邻居聚合的独立路径：局部路径（Local Pathway）和 DFS 路径。局部路径使用标准邻接矩阵 $A$ 处理直接邻居和即时连接；DFS 路径则通过重要性引导的 DFS 策略处理远端节点，该策略根据节点中心性度量对图拓扑进行有策略性的探索。

DFS 路径采用复杂的多步探索策略，根据节点的拓扑重要性对节点进行优先级排序。我们首先为每个节点计算综合重要性分数（Composite Importance Score）：
\begin{equation}
s_i = d_i + \beta \cdot b_i
\label{eq:importance_score}
\end{equation}
其中 $d_i = \sum_j A_{ij}$ 为度中心性，$b_i = \text{diag}(A^3)_i$ 作为介数中心性（Betweenness Centrality）的代理，$\beta = 0.3$ 为权重参数，用于平衡城市路网中局部连通性（度中心性）的主要重要性与战略位置（介数代理）的次要作用。

随后，我们使用统计阈值将节点划分为重要性层级：
\begin{equation}
\mathcal{H} = \{i : s_i > \mu_s + \sigma_s\}
\label{eq:high_importance}
\end{equation}
其中 $\mu_s$ 和 $\sigma_s$ 分别为重要性分数的均值和标准差，$\mathcal{H}$ 为高重要性节点集合。

多步 DFS 探索按层级递进：
\begin{equation}
L^{(t+1)} = \begin{cases}
L^{(t)}A \odot M_{\text{local}} & \text{若 } t < 2 \\
L^{(t)}A \odot (1 + S) & \text{若 } t \geq 2
\end{cases}
\label{eq:dfs_exploration}
\end{equation}
其中 $L^{(t)}$ 表示步骤 $t$ 时的可达性矩阵，$M_{\text{local}} = \text{diag}(1-\mathbb{I}_{\mathcal{H}})$ 在早期步骤中过滤局部连接，$S = s\mathbf{1}^T$ 为重要性分数矩阵，$\odot$ 表示逐元素乘法。

累积 DFS 邻接矩阵构建如下：
\begin{equation}
A_{\text{DFS}} = \sum_{t=1}^T L^{(t)}
\label{eq:dfs_adjacency}
\end{equation}
其中 $T$ 为最大 DFS 步数（默认 $T = 5$）。

对于每条路径，基于连接节点之间的双曲距离计算注意力系数：
\begin{equation}
\alpha_{ij}^{\text{local}} = \frac{\exp(-\sqrt{(d_{\mathcal{H}}(h_i,h_j))^2 + \epsilon})}{\sum_{k \in \mathcal{N}_{\text{local}}(i)}\exp(-\sqrt{(d_{\mathcal{H}}(h_i,h_k))^2 + \epsilon})}
\label{eq:local_attention}
\end{equation}
\begin{equation}
\alpha_{ij}^{\text{DFS}} = \frac{\exp(-\sqrt{(d_{\mathcal{H}}(h_i,h_j) \cdot w_{ij})^2 + \epsilon})}{\sum_{k \in \mathcal{N}_{\text{DFS}}(i)}\exp(-\sqrt{(d_{\mathcal{H}}(h_i,h_k) \cdot w_{ik})^2 + \epsilon})}
\label{eq:dfs_attention}
\end{equation}
其中 $d_{\mathcal{H}}(h_i,h_j)$ 表示节点嵌入之间的双曲距离，$\mathcal{N}_{\text{local}}(i)$ 和 $\mathcal{N}_{\text{DFS}}(i)$ 分别表示局部和 DFS 邻居集合，$w_{ij}$ 为动态调整权重，计算如下：
\begin{equation}
w_{ij} = (1 - \gamma \cdot \sigma(W_{\text{adj}}\bar{h}_i))
\label{eq:dynamic_weight}
\end{equation}
其中 $\gamma$ 为可学习的 DFS 权重参数，$\sigma$ 为 sigmoid 函数，$W_{\text{adj}}$ 为可学习变换矩阵，$\bar{h}_i$ 表示节点 $i$ 的平均特征向量。

数值稳定常数 $\epsilon = 10^{-8}$ 用于防止注意力计算中的数值不稳定性。这种重要性引导的 DFS 策略使模型能够通过对拓扑最重要节点的选择性探索，在保持计算效率的同时捕获局部结构模式和长程依赖关系。

\subsubsection{多层聚合}

遵循 GIN 范式 \cite{Xu2019}，我们引入可学习的自环重要性以增强模型的表达能力。GIN 通过确保聚合函数能够区分不同的邻居特征多重集（Multiset），解决了传统 GNN 的根本局限，这对于捕获路网中多样的连通模式至关重要。核心思想在于：添加可学习的自环权重允许模型在聚合邻域特征时保留中心节点的信息，避免因局部结构不同而节点产生相同表示所导致的信息损失。

\begin{equation}
h_i^{\text{self}} = W_{\text{self}}x_i
\label{eq:self_loop}
\end{equation}
\begin{equation}
h_i = h_i^{\text{neigh}} + (1+\epsilon)h_i^{\text{self}}
\label{eq:aggregation}
\end{equation}
其中 $W_{\text{self}}$ 为处理节点自身特征 $x_i$ 以生成自表示 $h_i^{\text{self}}$ 的可学习变换矩阵；$\epsilon$ 为可学习参数（初始化为 0.5），控制节点自身信息与邻域信息的相对重要性；$h_i^{\text{neigh}}$ 表示来自双路径注意力机制的聚合邻居信息。我们采用两层多层感知机（Multi-layer Perceptron，MLP）\cite{Goodfellow2016} 进行特征变换。

\subsubsection{自适应融合策略}

为充分利用来自局部和 DFS 路径的邻居信息，本文提出自适应融合机制以动态执行邻居聚合。对于每个注意力头 $k$，基于局部特征计算融合权重：
\begin{equation}
f^{(k)} = \text{mean}(H^{(k,\text{local})}, \text{axis}=-1)
\label{eq:fusion_feature}
\end{equation}
\begin{equation}
w^{(k)} = \sigma(W_{\text{fusion}}^{(k)} \cdot f^{(k)} + b_{\text{fusion}}^{(k)})
\label{eq:fusion_weight}
\end{equation}
其中 $W_{\text{fusion}}^{(k)} \in \mathbb{R}^{1 \times 1}$ 和 $b_{\text{fusion}}^{(k)}$ 为可学习参数，$\sigma$ 为 sigmoid 函数。

融合权重随后被扩展并逐元素应用：
\begin{equation}
w^{(k)} = \text{tile}(w^{(k)}, [1,1,D])
\label{eq:weight_expansion}
\end{equation}
\begin{equation}
H^{(k)} = w^{(k)} \odot H^{(k,\text{local})} + (1-w^{(k)}) \odot H^{(k,\text{DFS})}
\label{eq:adaptive_fusion}
\end{equation}
其中 $\odot$ 表示逐元素乘法。

\subsection{任务模块与模型训练}

HRGNN 采用专用任务模块（如分类层），旨在有效利用双曲空间中学到的层级表示，同时完成预测任务（如多类道路类型预测）。为弥合双曲表示与标准架构之间的差距，我们使用对数映射（Logarithmic Map）将最终双曲节点嵌入投影到原点处的切空间。该投影在保留双曲空间所捕获层级信息的同时，使表示与传统神经网络层兼容。

\subsubsection{作为任务模块的分类头}

以多类道路类型预测为例，阐述任务模块的设计。投影后的特征通过两层 MLP 组成的分类头处理。第一层应用 ReLU 激活函数引入非线性，随后进行批归一化（Batch Normalization）以稳定训练并提升收敛速度。Dropout 正则化用于防止过拟合，这对于路网中复杂层级模式的学习尤为重要。

第二层将隐含表示映射至对应不同道路类型的最终类别 logits。最后注意力层的双曲嵌入首先使用对数映射投影至欧几里得空间：
\begin{equation}
h_i^{\text{eucl}} = \log_0^c(h_i^H) = \frac{\tanh^{-1}(\sqrt{c}|h_i^H|)}{\sqrt{c}|h_i^H|}h_i^H
\label{eq:log_map}
\end{equation}

分类头由两层 MLP 组成：
\begin{equation}
z_i^{(1)} = \text{ReLU}(\text{BN}(W_1 h_i^{\text{eucl}} + b_1))
\label{eq:mlp_layer1}
\end{equation}
\begin{equation}
Z_i^{(1)} = \text{Dropout}(z_i^{(1)}, p=0.1)
\label{eq:dropout}
\end{equation}
\begin{equation}
\text{logits}_i = W_2 Z_i^{(1)} + b_2
\label{eq:mlp_layer2}
\end{equation}
其中 $W_1$、$W_2$ 为可学习权重矩阵，BN 表示批归一化，最终 logits 用于多类分类。

这一设计选择使模型能够学习复杂的决策边界，同时保持计算效率。MLP 架构轻量且具有足够的表达能力，能够基于层级位置和局部连通模式捕获不同道路类型之间的细微差别。

\subsubsection{模型训练}

模型优化采用复合损失函数，将用于分类的交叉熵损失与正则化项相结合。交叉熵损失为道路类型分类提供主要监督信号，L2 正则化防止参数过拟合。此外，本文引入曲率正则化，以鼓励模型学习能够反映路网层级本质的有意义双曲曲率值。

复合损失函数定义为
\begin{equation}
\mathcal{L} = \mathcal{L}_{\text{CE}} + \lambda_{\text{L2}}\mathcal{L}_{\text{L2}} + \lambda_{\text{curv}}\mathcal{L}_{\text{curv}}
\label{eq:total_loss}
\end{equation}
其中交叉熵损失为
\begin{equation}
\mathcal{L}_{\text{CE}} = -\frac{1}{|\mathcal{M}|}\sum_{i \in \mathcal{M}}\sum_{c=1}^C y_{ic}\log(\text{softmax}(\text{logits}_i)_c)
\label{eq:ce_loss}
\end{equation}
L2 正则化项为
\begin{equation}
\mathcal{L}_{\text{L2}} = \sum_{\theta \in \Theta}|\theta|_2^2
\label{eq:l2_loss}
\end{equation}
曲率正则化项为
\begin{equation}
\mathcal{L}_{\text{curv}} = (c - 1.0)^2
\label{eq:curv_loss}
\end{equation}
其中 $\mathcal{M}$ 表示用于训练的掩码节点集合，$C$ 为类别数，$y_{ic}$ 为真实标签，$\Theta$ 表示除曲率外的所有可训练参数，$\lambda_{\text{L2}}$、$\lambda_{\text{curv}}$ 为正则化系数。

双曲邻居聚合计算旨在最小化复合损失 $\mathcal{L}$，同时学习反映路网层级性质的最优曲率参数 $c$。学习目标是在满足双曲几何约束的同时，找到使道路类型预测分类精度最大化的参数。

训练过程在使用 Adam 优化器更新标准神经网络参数与在适当约束下更新双曲曲率参数之间交替进行，以确保数值稳定性。这种交替优化策略使模型能够通过曲率自适应联合学习层级结构，同时通过标准梯度下降学习分类决策边界。

\section{实验与结果}
\label{sec:experiments}

\subsection{实验任务}

实验评估针对双曲几何在路网表示学习有效性方面的两个核心研究问题：1）双曲方法是否在单个城市内相比欧几里得方法提供可量化的优势；2）双曲表示是否在具有不同形态特征的城市间表现出更优越的可迁移性。为回答上述问题，本文设计了两个互补的实验场景，分别考察局部表示学习能力和跨域泛化能力：通过城市内评估来检验局部层级模式的捕获能力，以及通过跨城市评估来检验向具有不同规划历史和地理约束的未见城市环境迁移的能力。

\begin{itemize}
\item \textbf{城市内道路类型分类}设置中，模型在同一图的不同节点子集上进行训练和评估，通过对节点设置训练、验证和测试标志来实现。

\item \textbf{跨城市道路类型分类}设置中，模型在特定城市的路网上训练，并在全新未见城市上进行评估。意大利开放街道地图（OpenStreetMap，OSM）数据集通过将 17 个意大利城市划分为训练、验证和测试集来演示这一方法，每个城市的节点均相应标注。这测试了模型向可能具有不同结构特性的路网泛化的能力。
\end{itemize}

\subsection{数据收集}

为确保数据集覆盖多样的几何特性，本文基于第三章所述的 $\delta$-双曲性分析结果选择代表性城市。分析揭示网络规模与 $\delta$-双曲性之间不存在强相关性，城市形态特征独立于网络规模，这为选择涵盖高、中、低双曲性的代表性城市提供了依据。

本文从开放街道地图（OSM）提取路网数据，该数据集提供了包含道路类型、限速和几何属性等丰富语义信息的全面地理数据 \cite{Boeing2017}。本研究汇编了具有不同特征的数据集，以评估模型在两种不同实验范式下的通用性。

表~\ref{tab:road_networks} 的可视化展示了所选城市中城市形态和层级路网结构的显著多样性，每个城市代表路网表示学习的独特挑战与机遇，共同验证了我们模型的有效性。

\begin{longtable}{@{\extracolsep{\fill}}ccc@{}}
\caption{城市内和跨城市数据集的路网结构。网络 (a)-(c) 代表具有不同 $\delta$-双曲性的城市内数据集；网络 (d)-(t) 展示具有多样层级模式的跨城市意大利数据集。}
\label{tab:road_networks} \\
\toprule
\multicolumn{3}{c}{\textbf{城市内数据集}} \\
\midrule
\endfirsthead
\multicolumn{3}{c}{{\tablename\ \thetable{} -- 续上页}} \\
\toprule
\endhead
\midrule
\multicolumn{3}{r}{{接下页}} \\
\endfoot
\bottomrule
\endlastfoot
% 城市内数据集
\includegraphics[width=0.22\textwidth]{figures/road/venice_road_network.png} &
\includegraphics[width=0.22\textwidth]{figures/road/munich_road_network.png} &
\includegraphics[width=0.22\textwidth]{figures/road/athens_road_network.png} \\[-5pt]
(a) 威尼斯 ($\delta=3.0$) & (b) 慕尼黑 ($\delta=17.5$) & (c) 雅典 ($\delta=32.0$) \\
\addlinespace[0.2em]
\midrule
\multicolumn{3}{c}{\textbf{跨城市数据集}} \\
\midrule
\addlinespace[0.1em]
% 跨城市数据集
\includegraphics[width=0.22\textwidth]{figures/road/italy_Orvieto_road_network.png} &
\includegraphics[width=0.22\textwidth]{figures/road/italy_Assisi_road_network.png} &
\includegraphics[width=0.22\textwidth]{figures/road/italy_San Gimignano_road_network.png} \\[-5pt]
(d) 奥尔维耶托 ($\delta=7.0$) & (e) 阿西西 ($\delta=7.0$) & (f) 圣吉米尼亚诺 ($\delta=8.0$) \\
\addlinespace[0.1em]
%
\includegraphics[width=0.22\textwidth]{figures/road/italy_Urbino_road_network.png} &
\includegraphics[width=0.22\textwidth]{figures/road/italy_Bergamo_road_network.png} &
\includegraphics[width=0.22\textwidth]{figures/road/italy_Perugia_road_network.png} \\[-5pt]
(g) 乌尔比诺 ($\delta=9.5$) & (h) 贝加莫 ($\delta=12.0$) & (i) 佩鲁贾 ($\delta=13.5$) \\
\addlinespace[0.1em]
%
\includegraphics[width=0.22\textwidth]{figures/road/italy_Genova_road_network.png} &
\includegraphics[width=0.22\textwidth]{figures/road/italy_Napoli_road_network.png} &
\includegraphics[width=0.22\textwidth]{figures/road/italy_Matera_road_network.png} \\[-5pt]
(j) 热那亚 ($\delta=14.0$) & (k) 那不勒斯 ($\delta=14.0$) & (l) 马泰拉 ($\delta=14.0$) \\
\addlinespace[0.1em]
%
\includegraphics[width=0.22\textwidth]{figures/road/italy_Pisa_road_network.png} &
\includegraphics[width=0.22\textwidth]{figures/road/italy_Verona_road_network.png} &
\includegraphics[width=0.22\textwidth]{figures/road/italy_Siena_road_network.png} \\[-5pt]
(m) 比萨 ($\delta=14.5$) & (n) 维罗纳 ($\delta=15.0$) & (o) 锡耶纳 ($\delta=16.0$) \\
\addlinespace[0.1em]
%
\includegraphics[width=0.22\textwidth]{figures/road/italy_Lucca_road_network.png} &
\includegraphics[width=0.22\textwidth]{figures/road/italy_Bologna_road_network.png} &
\includegraphics[width=0.22\textwidth]{figures/road/italy_Ravenna_road_network.png} \\[-5pt]
(p) 卢卡 ($\delta=16.0$) & (q) 博洛尼亚 ($\delta=18.0$) & (r) 拉文纳 ($\delta=19.5$) \\
\addlinespace[0.1em]
%
\includegraphics[width=0.22\textwidth]{figures/road/italy_Roma_road_network.png} &
\includegraphics[width=0.22\textwidth]{figures/road/italy_Firenze_road_network.png} &
\\[-5pt]
(s) 罗马 ($\delta=21.0$) & (t) 佛罗伦萨 ($\delta=27.5$) & \\
\end{longtable}

各数据集的特征统计见表~\ref{tab:dataset_stats}。

为理解数据集中路网的结构特性，本文分析了若干关键特征：

\begin{itemize}
\item \textbf{度分布：} 两个数据集均呈现重尾度分布（Heavy-tailed Degree Distribution）\cite{Barthelemy2011}，特征为少数高度节点（主要交叉口）与大量低度节点（断头路或简单交叉口）共存，近似遵循幂律分布。这一特性表明双曲表示可能特别适用。

\item \textbf{层级结构：} 路网中观察到清晰的层级模式 \cite{Han2020}，主要道路构成网络骨干，地方街道以树状方式围绕其布置。这种层级组织与双曲空间的指数展开特性高度契合。

\item \textbf{城市形态：} 数据集中的多个城市呈现出多样的城市形态 \cite{Zhang2023}，从具有有机不规则布局的历史中心到更具系统性规划的新区。这种多样性为测试模型适应不同结构模式的能力提供了具有挑战性的试验场。
\end{itemize}

\begin{table}[H]
\centering
\caption{路网数据集统计信息。}
\label{tab:dataset_stats}
\footnotesize
\begin{tabular}{@{}lcccc@{}}
\toprule
\textbf{统计项} & \textbf{威尼斯} & \textbf{慕尼黑} & \textbf{雅典} & \textbf{意大利城市} \\
\midrule
城市数量 & 1 & 1 & 1 & 17 \\
节点数 & 3,166 & 106,454 & 62,809 & 91,063 \\
边数 & 6,095 & 235,004 & 147,404 & 182,989 \\
每节点特征数 & 57 & 59 & 59 & 60 \\
道路类型数 & 6 & 12 & 12 & 5 \\
训练节点数 & 1,966 & 31,938 & 18,843 & 62,759 \\
验证节点数 & 400 & 37,258 & 21,983 & 6,197 \\
测试节点数 & 800 & 37,258 & 21,983 & 22,107 \\
平均度 & 3.85 & 4.42 & 4.69 & 4.02 \\
\bottomrule
\end{tabular}
\end{table}




\subsection{评估设置}

本文将 HRGNN 与涵盖传统 GNN、双曲方法和路网专用架构的全面基线方法集合进行对比，以确保严格的性能评估。

\begin{enumerate}
\item \textbf{基础 GNN。} 对比中包含四种代表图表示学习演进历程的基础 GNN 架构。图卷积网络（GCN）作为谱域方法的基础，利用归一化图拉普拉斯进行邻居聚合 \cite{Kipf2017}。GraphSAGE 提供通过采样和聚合局部邻居特征生成节点嵌入的归纳框架 \cite{Hamilton2017}。图注意力网络（GAT）在聚合过程中引入注意力机制，动态加权邻居重要性 \cite{Velickovic2018}。图同构网络（GIN）通过 MLP 建模单射多重集函数，为邻居聚合提供理论保障 \cite{Xu2019}。

\item \textbf{双曲神经网络。} 为专门评估双曲几何的贡献，本文引入两种先进的双曲神经网络。双曲注意力网络（HAT）在双曲空间中利用注意力机制进行层级表示学习 \cite{Gulcehre2019}。双曲图卷积网络（HGCN）将传统图卷积扩展至双曲空间，从而提升层级结构表示能力 \cite{Chami2019}。

\item \textbf{层级模型。} 考虑到交通网络的专业特性，本文引入四种专为路网分析设计的方法。关系融合网络（RFN）通过结合原始图和对偶图表示来建模路网 \cite{Jepsen2022}。GAIN 提供基于图注意力的通用 GNN 聚合机制改进 \cite{gharaee2021pr}。层级路网表示（HRNR）构建三层神经架构 \cite{Wu2020}。HyperRoad 采用双图模型，通过超图构建捕获高阶和长距离道路关系 \cite{Zhang2023}。

\item \textbf{简单基线模型。} 为建立性能界限并评估图学习的贡献，本文引入两种基本基线。随机基线（Random Baseline）分配随机类标签以评估性能下界；原始特征仅（Raw Features Only）直接将输入特征用于逻辑回归分类器，不进行任何图表示学习。
\end{enumerate}

所有基线方法均采用相同的实验协议和超参数优化程序以确保公平比较。实验协议采用穷举式超参数优化：系统探索学习率 $\{10^{-4}, 10^{-3}, 10^{-2}\}$；输出嵌入维度在 $\{64, 128, 256\}$ 之间评估，以平衡表达能力与计算效率；所有模型训练 2000 个轮次，Dropout 率统一为 0.1；批大小配置为城市内任务使用 1024，跨城市场景使用 2048。

对于 HRGNN，工程化的节点特征同时捕获路段的拓扑和几何属性，包括标准化几何向量、长度、中点坐标和限速。所有特征使用训练集的统计数据进行标准化。

性能评估采用三种互补指标，对分类能力进行全面评估：

\begin{itemize}
\item \textbf{微平均 F1 分数（Micro-averaged F1-Score）} 提供整体分类性能度量，对每个实例赋予相等权重，与类别无关。

\item \textbf{宏平均 F1 分数（Macro-averaged F1-Score）} 对每个类别赋予相等权重，与频率无关，鉴于道路类型之间存在显著的频率差异，该指标尤为重要。

\item \textbf{加权 F1 分数（Weighted F1-Score）} 计算按频率加权的类别 F1 分数，在微平均和宏平均指标之间提供平衡评估。
\end{itemize}

通过报告每个指标和数据集组合在 10 次不同随机权重初始化运行下的平均性能来确保统计稳健性。

\subsection{性能分析}

性能比较表中报告的所有结果均为 10 次不同随机权重初始化运行的平均性能，以确保统计稳健性。对于每个指标和数据集，最佳模型以粗体标注。如表~\ref{tab:performance_a} 和表~\ref{tab:performance_b}所示，HRGNN 在所有数据集和指标上持续取得最高分数。

% 表一：城市内分类 —— 威尼斯 & 慕尼黑
\begin{table}[H]
\centering
\caption{图神经网络模型在地理空间数据集上的性能比较（一）：城市内分类（威尼斯与慕尼黑）。最佳结果加粗，次佳结果加下划线，第三名以~* 标注。}
\label{tab:performance_a}
\begin{tabular}{@{}lcccccc@{}}
\toprule
& \multicolumn{6}{c}{\textbf{城市内分类}} \\
\cmidrule(lr){2-7}
& \multicolumn{3}{c}{\textbf{威尼斯}} & \multicolumn{3}{c}{\textbf{慕尼黑}} \\
\cmidrule(lr){2-4} \cmidrule(lr){5-7}
\textbf{模型} & 微平均 & 宏平均 & 加权 & 微平均 & 宏平均 & 加权 \\
& F1 & F1 & F1 & F1 & F1 & F1 \\
\midrule
随机基线 & 0.167 & 0.083 & 0.139 & 0.111 & 0.056 & 0.089 \\
原始特征 & 0.428 & 0.214 & 0.392 & 0.356 & 0.178 & 0.324 \\
\midrule
GCN      & 0.659             & 0.319             & 0.564             & 0.541             & 0.235             & 0.452 \\
GAT      & 0.699             & 0.408             & 0.684             & 0.597*            & \underline{0.342} & \underline{0.577} \\
GraphSAGE & 0.726*           & \underline{0.436} & \underline{0.716} & 0.589             & 0.367             & 0.574* \\
HAT      & 0.686             & 0.382             & 0.652             & 0.591             & 0.195             & 0.465 \\
HGCN     & 0.646             & 0.297             & 0.575             & 0.545             & 0.233             & 0.443 \\
RFN      & 0.616             & 0.218             & 0.520             & 0.587             & 0.190             & 0.456 \\
GIN      & 0.723             & \underline{0.436} & 0.709*            & 0.556             & 0.290             & 0.508 \\
GAIN     & 0.721             & 0.427*            & 0.700             & \underline{0.603} & 0.270             & 0.515 \\
HRNR     & 0.699             & 0.414             & 0.700             & 0.431             & 0.213             & 0.374 \\
HyperRoad & \underline{0.728} & 0.339            & 0.690             & 0.573             & 0.223             & 0.531 \\
\midrule
\textbf{HRGNN（本文）} & \textbf{0.786} & \textbf{0.477} & \textbf{0.779} & \textbf{0.658} & \textbf{0.336}* & \textbf{0.618} \\
\bottomrule
\end{tabular}
\end{table}


% 表二：城市内分类 —— 雅典 & 跨城市分类 —— 意大利城市
\begin{table}[H]
\centering
\caption{图神经网络模型在地理空间数据集上的性能比较（二）：城市内分类（雅典）与跨城市分类（意大利城市）。最佳结果加粗，次佳结果加下划线，第三名以~* 标注。}
\label{tab:performance_b}
\begin{tabular}{@{}lcccccc@{}}
\toprule
& \multicolumn{3}{c}{\textbf{城市内分类}} & \multicolumn{3}{c}{\textbf{跨城市分类}} \\
\cmidrule(lr){2-4} \cmidrule(lr){5-7}
& \multicolumn{3}{c}{\textbf{雅典}} & \multicolumn{3}{c}{\textbf{意大利城市}} \\
\cmidrule(lr){2-4} \cmidrule(lr){5-7}
\textbf{模型} & 微平均 & 宏平均 & 加权 & 微平均 & 宏平均 & 加权 \\
& F1 & F1 & F1 & F1 & F1 & F1 \\
\midrule
随机基线 & 0.125 & 0.062 & 0.098 & 0.162 & 0.043 & 0.120 \\
原始特征 & 0.312 & 0.156 & 0.285 & 0.283 & 0.142 & 0.255 \\
\midrule
GCN       & 0.553             & 0.133             & 0.470             & 0.421             & 0.260             & 0.414 \\
GAT       & 0.567             & 0.225             & 0.539             & 0.505             & 0.283             & 0.494 \\
GraphSAGE & 0.583             & 0.255             & 0.563             & 0.489             & 0.291             & 0.496 \\
HAT       & \underline{0.609} & 0.170             & 0.541             & \underline{0.565} & 0.261             & 0.515 \\
HGCN      & 0.591             & 0.175             & 0.527             & 0.465             & 0.261             & 0.461 \\
RFN       & 0.571             & 0.156             & 0.504             & 0.524             & 0.155             & 0.402 \\
GIN       & 0.606*            & \underline{0.274} & 0.571             & 0.560             & 0.306             & 0.540 \\
GAIN      & 0.605             & 0.277             & 0.573*            & 0.552             & \underline{0.299} & 0.527* \\
HRNR      & 0.569             & 0.180             & 0.521             & 0.562*            & 0.144             & 0.404 \\
HyperRoad & 0.578             & 0.263*            & \underline{0.577} & 0.555             & 0.181             & 0.440 \\
\midrule
\textbf{HRGNN（本文）} & \textbf{0.660} & \underline{\textbf{0.274}} & \textbf{0.624} & \textbf{0.575} & \textbf{0.295}* & \underline{\textbf{0.534}} \\
\bottomrule
\end{tabular}
\end{table}

HRGNN 在所有三个城市内分类数据集上均展现出卓越性能。在威尼斯数据集上，HRGNN 实现 0.786 的微平均 F1 分数，分别超越最佳传统方法 GraphSAGE（0.726）8.4\%、GCN（0.659）19.3\%，以及双曲方法 HAT（0.686）和 HGCN（0.646）14.6\% 和 21.7\%。

在慕尼黑数据集上，HRGNN 达到 0.658 的微平均 F1，超越 GAT（0.597）10.2\% 和 GraphSAGE（0.589）11.7\%。在雅典数据集上，HRGNN 实现 0.660 的微平均 F1，超越最佳传统方法 HAT（0.609）8.4\%，并分别超越其他强基线 GIN（0.606）和 GAIN（0.605）8.9\% 和 9.1\%。

在具有挑战性的意大利城市迁移学习场景中，HRGNN 实现 0.575 的微平均 F1，超越大多数基线。对比双曲方法揭示了一个关键发现：HRGNN 在微平均 F1 上超越 HAT 1.8\%（0.575 对比 0.565），在宏平均 F1 上取得 13.0\% 的提升。

HRGNN 在多样城市形态上的持续优越性，可归因于与路网内在特性相契合的若干基本理论原理。相比欧几里得基线 8.4\%--9.1\% 的性能提升（城市内数据集），直接反映了双曲空间在无失真表示树状结构方面相比传统欧几里得嵌入的理论优势。

城市路网的 $\delta$-双曲性与双曲方法的性能增益相关。呈现更强树状特征的城市（威尼斯：$\delta=3.0$；锡耶纳：$\delta=14.0$）相比网格状城市（雅典：$\delta=32.0$）展现出更大的改进幅度。这源于双曲空间指数展开特性与层级道路系统分支模式的匹配。此外，路网呈现幂律度分布——少数高连通度交叉口与大量低度节点共存。双曲嵌入通过允许中心节点在保持层级组织的同时与众多邻居维持合理距离，来适应这种异质性。

双路径架构整合了互补信息：局部路径捕获用于道路类型区分的即时连通模式，而 DFS 路径揭示用于基于层级分类的结构上下文。DFS 相比其他采样方法表现更优，原因在于其定向探索遵循城市网络中的自然流动模式，不同于对所有邻居一视同仁的广度优先扩展。这在需要向未见城市形态泛化的迁移学习场景中尤为有价值。重要性评分机制识别出构成城市出行骨干的关键连接器节点（对应主要交叉口和干道交汇点）。

\begingroup
\footnotesize
\setlength{\tabcolsep}{2pt}
\renewcommand{\arraystretch}{1.2}
\begin{longtable}{|>{\centering\arraybackslash}m{1.45cm}|>{\centering\arraybackslash}m{2cm}|>{\centering\arraybackslash}m{2cm}|>{\centering\arraybackslash}m{2cm}|>{\centering\arraybackslash}m{2cm}|>{\centering\arraybackslash}m{2cm}|}
\caption{各模型在测试城市样本道路上的预测结果。}
\label{tab:model_comparison} \\
\hline
\multirow{2}{*}{\textbf{模型}} & \multicolumn{3}{c|}{\textbf{城市内分类}} & \multicolumn{2}{c|}{\textbf{跨城市分类}} \\
\cline{2-6}
& \textbf{威尼斯} & \textbf{慕尼黑} & \textbf{雅典} & \textbf{罗马} & \textbf{热那亚} \\
& \textbf{样本} & \textbf{样本} & \textbf{样本} & \textbf{样本} & \textbf{样本} \\
\hline
\endfirsthead

\multicolumn{6}{c}%
{{\bfseries \tablename\ \thetable{} -- 续上页}} \\
\hline
\multirow{2}{*}{\textbf{模型}} & \multicolumn{3}{c|}{\textbf{城市内分类}} & \multicolumn{2}{c|}{\textbf{跨城市分类}} \\
\cline{2-6}
& \textbf{威尼斯} & \textbf{慕尼黑} & \textbf{雅典} & \textbf{罗马} & \textbf{热那亚} \\
& \textbf{样本} & \textbf{样本} & \textbf{样本} & \textbf{样本} & \textbf{样本} \\
\hline
\endhead

\hline \multicolumn{6}{|r|}{{接下页}} \\ \hline
\endfoot

\hline
\multicolumn{6}{l}{%
\makebox[\textwidth][s]{%
\textcolor{green}{\rule{0.4cm}{0.12cm}} 训练集 \hfill
\textcolor{blue}{\rule{0.4cm}{0.12cm}} 验证集 \hfill
\textcolor{yellow}{\rule{0.4cm}{0.12cm}} 预测正确 \hfill
\textcolor{red}{\rule{0.4cm}{0.12cm}} 预测错误%
}%
} \\
\endlastfoot

\textbf{真实值} & 
\includegraphics[width=0.15\textwidth]{figures/model/1.png} & 
\includegraphics[width=0.15\textwidth]{figures/model/2.png} & 
\includegraphics[width=0.15\textwidth]{figures/model/3.png} & 
\includegraphics[width=0.15\textwidth]{figures/model/4.png} & 
\includegraphics[width=0.15\textwidth]{figures/model/5.png} \\
\hline

\textbf{GCN} & 
\includegraphics[width=0.15\textwidth]{figures/model/GCN1.png} & 
\includegraphics[width=0.15\textwidth]{figures/model/GCN2.png} & 
\includegraphics[width=0.15\textwidth]{figures/model/GCN3.png} & 
\includegraphics[width=0.15\textwidth]{figures/model/GCN4.png} & 
\includegraphics[width=0.15\textwidth]{figures/model/GCN5.png} \\
\hline

\textbf{\footnotesize GraphSAGE} & 
\includegraphics[width=0.15\textwidth]{figures/model/GraphSAGE1.png} & 
\includegraphics[width=0.15\textwidth]{figures/model/GraphSAGE2.png} & 
\includegraphics[width=0.15\textwidth]{figures/model/GraphSAGE3.png} & 
\includegraphics[width=0.15\textwidth]{figures/model/GraphSAGE4.png} & 
\includegraphics[width=0.15\textwidth]{figures/model/GraphSAGE5.png} \\
\hline

\textbf{HGCN} & 
\includegraphics[width=0.15\textwidth]{figures/model/HGCN1.png} & 
\includegraphics[width=0.15\textwidth]{figures/model/HGCN2.png} & 
\includegraphics[width=0.15\textwidth]{figures/model/HGCN3.png} & 
\includegraphics[width=0.15\textwidth]{figures/model/HGCN4.png} & 
\includegraphics[width=0.15\textwidth]{figures/model/HGCN5.png} \\
\hline

\textbf{RFN} & 
\includegraphics[width=0.15\textwidth]{figures/model/RFN1.png} & 
\includegraphics[width=0.15\textwidth]{figures/model/RFN2.png} & 
\includegraphics[width=0.15\textwidth]{figures/model/RFN3.png} & 
\includegraphics[width=0.15\textwidth]{figures/model/RFN4.png} & 
\includegraphics[width=0.15\textwidth]{figures/model/RFN5.png} \\
\hline

\textbf{GAIN} & 
\includegraphics[width=0.15\textwidth]{figures/model/GAIN1.png} & 
\includegraphics[width=0.15\textwidth]{figures/model/GAIN2.png} & 
\includegraphics[width=0.15\textwidth]{figures/model/GAIN3.png} & 
\includegraphics[width=0.15\textwidth]{figures/model/GAIN4.png} & 
\includegraphics[width=0.15\textwidth]{figures/model/GAIN5.png} \\
\hline

\textbf{HRNR} & 
\includegraphics[width=0.15\textwidth]{figures/model/HRNR1.png} & 
\includegraphics[width=0.15\textwidth]{figures/model/HRNR2.png} & 
\includegraphics[width=0.15\textwidth]{figures/model/HRNR3.png} & 
\includegraphics[width=0.15\textwidth]{figures/model/HRNR4.png} & 
\includegraphics[width=0.15\textwidth]{figures/model/HRNR5.png} \\
\hline

\textbf{\footnotesize HyperRoad} & 
\includegraphics[width=0.15\textwidth]{figures/model/HyperRoad1.png} & 
\includegraphics[width=0.15\textwidth]{figures/model/HyperRoad2.png} & 
\includegraphics[width=0.15\textwidth]{figures/model/HyperRoad3.png} & 
\includegraphics[width=0.15\textwidth]{figures/model/HyperRoad4.png} & 
\includegraphics[width=0.15\textwidth]{figures/model/HyperRoad5.png} \\
\hline

\textbf{HRGNN} & 
\includegraphics[width=0.15\textwidth]{figures/model/HRGNN1.png} & 
\includegraphics[width=0.15\textwidth]{figures/model/HRGNN2.png} & 
\includegraphics[width=0.15\textwidth]{figures/model/HRGNN3.png} & 
\includegraphics[width=0.15\textwidth]{figures/model/HRGNN4.png} & 
\includegraphics[width=0.15\textwidth]{figures/model/HRGNN5.png} \\
\hline

\end{longtable}
\endgroup

为补充定量评估，本文开展了案例研究，考察代表性城市的预测模式。表~\ref{tab:model_comparison} 呈现了八种模型在五个具有不同几何属性城市上的对比结果。

威尼斯（$\delta=3.0$）因其不规则中世纪布局而构成最具挑战性的案例。HRGNN 实现了主导性覆盖，仅在复杂交叉口处出现零散错误；而欧几里得基线（GCN、GraphSAGE、GAIN）在密集小巷区域呈现连续错误带。这反映了欧几里得多项式展开与指数层级分支之间的根本不匹配。尽管 HGCN 在双曲空间中运行，其错误密度与欧几里得基线相当，证实了仅靠空间变换而缺乏适当注意力机制是不足的。

慕尼黑（$\delta=17.5$）将放射状层级结构与网格状区域相结合。HRGNN 在两种形态上均保持均衡覆盖；而竞争方法显示出形态依赖的性能差异：HGCN 在放射状区域成功但在网格外围失败，HyperRoad 则表现相反。这验证了 HRGNN 自适应融合策略处理多样城市模式的有效性。

雅典（$\delta=32.0$）测试了双曲方法在低双曲性网络中是否会受到惩罚。尽管具有欧几里得特性，HRGNN 仍展现出广泛覆盖，验证了可学习曲率参数能够成功适应不同的网络特征。向罗马（$\delta=21.0$）和热那亚（$\delta=14.0$）的跨城市迁移进一步证明了 HRGNN 在网格状结构和地形约束不规则形态下的鲁棒性。

模型失败集中在复杂交叉口、层级边界和密集小巷。HRGNN 在真正模糊的情形下出现孤立错误，而竞争者则在特定结构模式上表现出空间聚集性失败。视觉性能差距与网络双曲性相关，HRGNN 在慕尼黑双重形态上的成功表明，来自局部和 DFS 路径的互补信息提供了单策略方法所不具备的鲁棒性。

\subsection{消融实验}

本文开展了全面的消融实验，以评估 HRGNN 各关键组件的独立贡献。表~\ref{tab:ablation} 呈现了系统性移除各主要组件以量化其对道路类型分类性能影响的结果。

\begin{table}[t]
\centering
\caption{消融实验结果——各变体性能比较（微平均 F1 分数）。}
\label{tab:ablation}
\footnotesize
\begin{tabular}{@{}lcccc@{}}
\toprule
\textbf{模型变体} & \textbf{威尼斯} & \textbf{慕尼黑} & \textbf{雅典} & \textbf{意大利城市} \\
\midrule
HRGNN（完整模型）& 0.78 & 0.65 & 0.63 & 0.57 \\
\midrule
去除双曲空间（欧几里得）& 0.73 & 0.59 & 0.58 & 0.53 \\
去除注意力机制 & 0.73 & 0.58 & 0.57 & 0.53 \\
去除 DFS 邻居 & 0.76 & 0.62 & 0.57 & 0.55 \\
\bottomrule
\end{tabular}
\end{table}

消融分析揭示了三个关键发现，共同验证了 HRGNN 的架构设计选择。当双曲操作被替换为欧几里得对应操作时，所有数据集上的性能一致下降 6.4\%--9.2\%，证实了双曲空间对层级结构的天然亲和性相比传统欧几里得嵌入为路网提供了实质性优势。类似地，移除注意力机制并施加统一聚合权重导致了相当幅度的性能下降（6.4\%--10.8\%），表明选择性信息聚合对于区分城市网络中不同道路连接的变化重要性至关重要。尽管 DFS 邻居组件的改进幅度相对较小但持续（2.6\%--4.6\%），它仍对捕获超越即时局部连通性的更广泛结构上下文做出了有意义的贡献。

\subsection{DFS 采样策略评估}

为验证双路径注意力机制的有效性，本文开展了系统性实验，以确定最优 DFS 配置，并在四个数据集上比较替代性全局路径采样方法。

\subsubsection{采样深度}

本文研究了采样深度（即 DFS 步数 $T$）对模型性能的影响，在城市内和跨城市场景中测试了 1 至 8 步的值。每个实验使用不同随机种子重复 5 次以确保统计可靠性。

\begin{table}[H]
\centering
\caption{DFS 步数对四个数据集性能的影响（微平均 F1 分数）。}
\label{tab:dfs_steps}
\footnotesize
\begin{tabular}{@{}ccccc@{}}
\toprule
\textbf{DFS 步数 ($T$)} & \textbf{威尼斯} & \textbf{慕尼黑} & \textbf{雅典} & \textbf{意大利城市} \\
\midrule
1 & 0.771 & 0.643 & 0.609 & 0.519 \\
2 & 0.778 & 0.651 & 0.616 & 0.527 \\
3 & 0.782 & 0.655 & 0.621 & 0.532 \\
4 & 0.785 & 0.657 & 0.623 & 0.534 \\
5 & \textbf{0.786} & \textbf{0.658} & \textbf{0.624} & \textbf{0.534} \\
6 & 0.784 & 0.656 & 0.622 & 0.532 \\
7 & 0.781 & 0.653 & 0.619 & 0.529 \\
8 & 0.779 & 0.650 & 0.617 & 0.526 \\
\bottomrule
\end{tabular}
\end{table}

如表~\ref{tab:dfs_steps} 所示，实验结果表明在所有四个数据集上，性能从 $T=1$ 到 $T=5$ 持续提升，并在 $T=5$ 时达到峰值。超过 $T=5$ 后，性能开始下降。这种在多样城市形态中一致的规律表明，虽然多步采样比单步方法捕获了更丰富的结构信息，但过度探索（$T>5$）会引入来自无关远端邻居的噪声。

\subsubsection{采样策略}

为进一步验证重要性引导 DFS 方法的有效性，本文在双路径框架内将其与五种替代性全局路径采样方法进行了比较。表~\ref{tab:sampling_strategy} 总结了性能比较结果。

\begin{table}[H]
\centering
\caption{四个数据集上全局路径采样方法比较（微平均 F1 分数）。}
\label{tab:sampling_strategy}
\footnotesize
\begin{tabular}{@{}lcccc@{}}
\toprule
\textbf{方法} & \textbf{威尼斯} & \textbf{慕尼黑} & \textbf{雅典} & \textbf{意大利城市} \\
\midrule
HRGNN-DFS（本文）& \textbf{0.786} & \textbf{0.658} & \textbf{0.624} & \textbf{0.534} \\
个性化 PageRank & 0.779 & 0.648 & 0.616 & 0.524 \\
广度优先采样 & 0.768 & 0.638 & 0.607 & 0.514 \\
基于度中心性 & 0.762 & 0.632 & 0.601 & 0.508 \\
带重启的随机游走 & 0.758 & 0.628 & 0.597 & 0.502 \\
随机多跳 & 0.745 & 0.615 & 0.585 & 0.489 \\
\bottomrule
\end{tabular}
\end{table}

对比分析表明，重要性引导 DFS 采样在四个数据集上持续超越所有替代方法。性能差距在跨城市迁移学习（意大利城市）中最为显著，DFS 保持了更大的优势，表明其对未见城市网络具有更优越的泛化能力。

上述发现确认 $T=5$ 为最优 DFS 配置，并验证了重要性引导 DFS 采样在多样城市形态的路网表示学习中优于传统图采样方法。


\section{本章小结}

本章提出了双曲路网图神经网络（HRGNN），通过庞加莱球模型中的双曲几何建模路网的层次结构，并设计了双通道双曲注意力机制分别捕获局部连通性和全局层次依赖。实验表明，隐式利用双曲空间的几何特性即可带来显著的性能提升，验证了双曲几何对路网层次结构建模的有效性。同时，本章分析也揭示了隐式范式在层次关系显式表达方面的不足，为第五章引入蕴含锥机制进行显式层次约束奠定了基础。